\chapter{Sécurité applicative et RGPD}

\section{Protection contre les vulnérabilités OWASP}
Pour sécuriser \textbf{Badge Reader}, j’ai appliqué les recommandations \textbf{OWASP Top 10} en utilisant les outils natifs d’Elixir/Phoenix et des bonnes pratiques spécifiques.
\newline

\begin{for_you}{Mes protections contre les vulnérabilités courantes}
    \textbf{Sécurité intégrée à Phoenix} :
    \begin{itemize}
        \item \textbf{XSS} : Échappement automatique des templates HEEx (équivalent à React).
        \item \textbf{CSRF} : Tokens intégrés dans les formulaires et requêtes.
        \item \textbf{SQL Injection} : Requêtes paramétrées avec Ecto (ORM d’Elixir).
        \item \textbf{Headers de sécurité} : Configuration via \texttt{Plug.Headers}.
    \end{itemize}
\end{for_you}

\begin{for_you}{Exemple d'éléments mie en place}
\textit{Protection XSS} :
    \begin{lstlisting}
    # Dans un template HEEx (lib/badge_reader_web/live/user_live.ex)
    def render(assigns) do
      ~H"""
      <div>
        <!-- Les données sont automatiquement échappées -->
        <h1><%= @user.name %></h1>
        <p><%= @user.bio %></p>

        <!-- Danger : utiliser raw/1 uniquement pour du HTML sûr -->
        <%= raw(@trusted_html) %>
      </div>
      """
    end
    \end{lstlisting}

    \textit{Protection CSRF} :
    \begin{lstlisting}
    # Dans le router.ex
    pipeline :browser do
      plug :accepts, ["html"]
      plug :fetch_session
      plug :fetch_flash
      plug :protect_from_forgery  # Protection CSRF intégrée
      plug :put_secure_browser_headers  # Headers de sécurité
    end
    \end{lstlisting}

    \textit{Requête Ecto (anti-SQL Injection)} :
    \begin{lstlisting}
    # Dans un repository (lib/badge_reader/repos/users.ex)
    def get_user_by_email(email) do
      BadgeReader.Repo.one(
        from u in BadgeReader.Accounts.User,
        where: u.email == ^email,  # ^ force le binding sécurisé
        select: u
      )
    end
    \end{lstlisting}

    \textbf{Headers de sécurité} (dans \texttt{config/config.exs}) :
    \begin{lstlisting}
    config :badge_reader, BadgeReaderWeb.Endpoint,
      url: [host: "localhost"],
      http: [
        port: 4000,
        # Headers de sécurité
        :secure_browser_headers => %{
          "x-frame-options" => "DENY",
          "x-content-type-options" => "nosniff",
          "x-xss-protection" => "1; mode=block",
          "content-security-policy" => "default-src 'self'"
        }
      ]
    \end{lstlisting}
\end{for_you}

\section{Authentification et autorisation}

L'authentification utilise JWT avec des tokens d'accès courts (15 minutes) et des refresh tokens sécurisés. Le hachage des mots de passe utilise Argon2, plus sécurisé que bcrypt. L'autorisation implémente un système de rôles et permissions granulaire avec des middlewares de vérification.

La gestion des sessions sécurise les tokens avec des cookies HttpOnly et SameSite. La déconnexion invalide les tokens côté serveur et côté client pour garantir la sécurité.

\begin{exemple}
\textbf{Configuration JWT et Argon2 (1/3) :}
\begin{lstlisting}[language=JavaScript]
const jwt = require('jsonwebtoken');
const argon2 = require('argon2');

// Configuration JWT
const JWT_SECRET = process.env.JWT_SECRET;
const JWT_EXPIRES_IN = '15m';
const REFRESH_EXPIRES_IN = '7d';

// Hachage des mots de passe avec Argon2
const hashPassword = async (password) => {
  return await argon2.hash(password, {
    type: argon2.argon2id,
    memoryCost: 2 ** 16, // 64 MB
    timeCost: 3,
    parallelism: 1
  });
};
\end{lstlisting}
\end{exemple}

\begin{exemple}
\textbf{Configuration JWT et Argon2 (2/3) :}
\begin{lstlisting}[language=JavaScript]
// Génération des tokens
const generateTokens = (userId, role) => {
  const accessToken = jwt.sign(
    { userId, role, type: 'access' },
    JWT_SECRET,
    { expiresIn: JWT_EXPIRES_IN }
  );

  const refreshToken = jwt.sign(
    { userId, type: 'refresh' },
    JWT_SECRET,
    { expiresIn: REFRESH_EXPIRES_IN }
  );

  return { accessToken, refreshToken };
};
\end{lstlisting}
\end{exemple}

\begin{exemple}
\textbf{Configuration JWT et Argon2 (3/3) :}
\begin{lstlisting}[language=JavaScript]
// Middleware d'authentification
const authenticateToken = (req, res, next) => {
  const authHeader = req.headers['authorization'];
  const token = authHeader && authHeader.split(' ')[1];

  if (!token) {
    return res.status(401).json({
      error: 'Token d\'accès requis'
    });
  }

  jwt.verify(token, JWT_SECRET, (err, user) => {
    if (err) {
      return res.status(403).json({
        error: 'Token invalide'
      });
    }
    req.user = user;
    next();
  });
};
\end{lstlisting}

\textbf{Système de permissions (1/2) :}
\begin{lstlisting}[language=JavaScript]
// Définition des permissions
const PERMISSIONS = {
  PROJECT_CREATE: 'project:create',
  PROJECT_READ: 'project:read',
  PROJECT_UPDATE: 'project:update',
  PROJECT_DELETE: 'project:delete',
  USER_MANAGE: 'user:manage'
};

// Rôles et leurs permissions
const ROLES = {
  ADMIN: [PERMISSIONS.PROJECT_CREATE, PERMISSIONS.PROJECT_READ,
          PERMISSIONS.PROJECT_UPDATE, PERMISSIONS.PROJECT_DELETE,
          PERMISSIONS.USER_MANAGE],
  MANAGER: [PERMISSIONS.PROJECT_CREATE, PERMISSIONS.PROJECT_READ,
            PERMISSIONS.PROJECT_UPDATE],
  DEVELOPER: [PERMISSIONS.PROJECT_READ, PERMISSIONS.PROJECT_UPDATE]
};
\end{lstlisting}

\textbf{Système de permissions (2/2) :}
\begin{lstlisting}[language=JavaScript]
// Middleware d'autorisation
const requirePermission = (permission) => {
  return (req, res, next) => {
    const userPermissions = ROLES[req.user.role] || [];
    if (!userPermissions.includes(permission)) {
      return res.status(403).json({
        error: 'Permissions insuffisantes'
      });
    }
    next();
  };
};
\end{lstlisting}
\end{exemple}

\begin{conseil}
\begin{itemize}
    \item Utiliser JWT avec des tokens courts et refresh tokens
    \item Implémenter Argon2 pour le hachage des mots de passe
    \item Créer un système de rôles et permissions granulaire
    \item Sécuriser les cookies avec HttpOnly et SameSite
    \item Implémenter la déconnexion sécurisée
\end{itemize}
\end{conseil}

\begin{jury}
\begin{itemize}
    \item Pourquoi utiliser JWT plutôt que des sessions ?
    \item Comment gérez-vous la sécurité des mots de passe ?
    \item Votre système d'autorisation est-il granulaire ?
    \item Comment gérez-vous l'expiration des tokens ?
    \item Avez-vous prévu la révocation des tokens ?
\end{itemize}
\end{jury}

\section{Conformité RGPD}

Pour respecter le \textbf{RGPD}, j’ai implémenté :
\begin{itemize}
    \item Un \textbf{registre des traitements}.
    \item La \textbf{minimisation des données}.
    \item Les \textbf{droits des utilisateurs} (accès, rectification, effacement).
    \item Le \textbf{chiffrement} des données sensibles.
    \newline
\end{itemize}

\begin{for_you}{1. Registre des traitements :}
    \begin{lstlisting}
    # Dans lib/badge_reader/privacy.ex
    defmodule BadgeReader.Privacy do
      @processing_activities %{
        "user_management" => %{
          purpose: "Gestion des comptes utilisateurs",
          legal_basis: "Consentement",
          data_categories: ["identité", "authentification"],
          retention: "3 ans après désactivation",
          recipients: ["administrateurs", "hébergeur"]
        },
        "access_logs" => %{
          purpose: "Traçabilité des accès",
          legal_basis: "Obligation légale",
          data_categories: ["logs", "horodatage"],
          retention: "1 an",
          recipients: ["administrateurs"]
        }
      }

      def get_processing_activity(activity_name) do
        @processing_activities[activity_name]
      end
    end
    \end{lstlisting}
\end{for_you}

\begin{for_you}{2. Implémentation des droits RGPD :}
    \begin{lstlisting}
    # Dans lib/badge_reader/accounts.ex
    defmodule BadgeReader.Accounts do
      def get_personal_data(user_id) do
        Repo.one(
          from u in User,
          where: u.id == ^user_id,
          select: %{email: u.email, role: u.role, created_at: u.inserted_at}
        )
      end

      def delete_personal_data(user_id) do
        Repo.transaction(fn ->
          Repo.delete_all(from b in Badge, where: b.user_id == ^user_id)
          Repo.delete!(User, id: user_id)
        end)
      end

      def export_personal_data(user_id) do
        user = Repo.get!(User, user_id)
        badges = Repo.all(from b in Badge, where: b.user_id == ^user_id)
        %{user: user, badges: badges}
      end
    end
    \end{lstlisting}
\end{for_you}

\textbf{\newline}

\begin{for_you}{3. Chiffrement des données sensibles :}
    \begin{lstlisting}
    # Dans lib/badge_reader/security.ex
    defmodule BadgeReader.Security do
      @aes_key Application.compile_env(:badge_reader, :aes_key)

      def encrypt(data) do
        {:ok, iv, ciphertext} = :crypto.block_encrypt(:aes_256_cbc, @aes_key, data)
        Base.encode64(iv <> ciphertext)
      end

      def decrypt(encoded) do
        decoded = Base.decode64!(encoded)
        <<iv::binary-16, ciphertext::binary>> = decoded
        :crypto.block_decrypt(:aes_256_cbc, @aes_key, iv, ciphertext)
      end
    end
    \end{lstlisting}
\end{for_you}

\begin{for_you}{4. Notification des violations :}
    \begin{lstlisting}
    # Dans lib/badge_reader/alerts.ex
    defmodule BadgeReader.Alerts do
      def notify_breach(breach_details) do
        # Envoi d'alerte par email (exemple simplifié)
        email = %{
          to: "dpo@colint.school",
          subject: "Violation de données - Badge Reader",
          body: "Une violation a été détectée : #{inspect(breach_details)}"
        }
        Mailer.deliver(email)
      end
    end
    \end{lstlisting}
\end{for_you}
\section{Liens utiles}

\begin{itemize}
    \item OWASP Top 10: \url{https://owasp.org/www-project-top-ten/}
    \item OWASP Cheat Sheets: \url{https://cheatsheetseries.owasp.org/}
    \item CNIL RGPD: \url{https://www.cnil.fr/fr/rgpd-de-quoi-parle-t-on}
    \item Argon2: \url{https://github.com/P-H-C/phc-winner-argon2}
    \item JWT Best Practices: \url{https://tools.ietf.org/html/rfc8725}
\end{itemize}
